\begin{abstract}
Biologically-inspired AI agent frameworks claim reliability benefits
through structural guarantees adapted from gene regulatory networks,
immune systems, and metabolic control.  These claims are rarely tested
empirically against simpler alternatives.  We present three deep
benchmarks---metabolic priority gating, autoinducer-based quorum
sensing, and Bayesian stagnation detection---each comparing a
biologically-grounded implementation against a naive non-biological
alternative and an ablated control, across 1{,}000 trials per seed
and 10 seeds (10M+ data points total).  Metabolic priority gating
delivers 100\% critical-operation service under bursty load versus
39.8\% for a flat budget counter ($\Delta = +0.602$, $p < 0.001$).
Quorum sensing achieves 0\% false-positive rate with 71--87\%
true-positive rate at 40\% agent compromise, occupying a unique
precision--recall operating point that neither majority voting nor
independent detection reaches.  Bayesian stagnation detection loses to cosine similarity with mock
embeddings but dominates with real sentence embeddings: 96\%
false-stagnation discrimination versus 2\% naive, and 96\% convergence
accuracy versus 40\% naive.  This reveals a conditional structural
guarantee that activates when embedding quality is present.  The
pattern across all three benchmarks: biological design earns its
complexity through \emph{mechanism-level structural guarantees}---priority
gating, signal accumulation with temporal decay, two-signal
discrimination---rather than through algorithmic sophistication.
\end{abstract}
